% ============================================================
% CHƯƠNG 2 — CÔNG NGHỆ SỬ DỤNG  (mẫu SoICT chương 3)
% Sắp xếp lại ngày 07/11/2026 theo mẫu Do_An_Latex_SoiCT. Bản trước: backup/2026-11-07_truoc-mau-soict/Chap2.tex
% (tên cũ "CƠ SỞ LÝ THUYẾT VÀ CÔNG NGHỆ"). Mục "Các nghiên cứu và hệ thống liên quan" (sec:lienquan:
% Vaccargiu, Liu, Toucan, Tsai) đã chuyển, rút gọn, vào Chương 1 mục 1.5.4 (subsec:verra-undp) vì mẫu SoICT
% đặt khảo sát hệ thống tương tự ở chương khảo sát; bản đầy đủ vẫn nằm trong file backup.
%
% Hướng dẫn của mẫu SoICT cho chương này:
%   - không quá 10 trang; nội dung dài hơn đưa vào Phụ lục; phân tích và tóm tắt, không trình bày dài dòng;
%   - chỉ giới thiệu công nghệ, nền tảng, thư viện mà đồ án thực sự dùng; không trình bày lại kiến thức
%     phổ thông hay sách giáo khoa;
%   - với mỗi công nghệ: (i) nó giải quyết vấn đề / yêu cầu nào ở Chương 1 (NC-xx, VĐ-x, yêu cầu phi chức năng),
%     (ii) các phương án thay thế, (iii) lý do chọn; có thể dùng bảng so sánh;
%   - phần thiết kế, cài đặt cụ thể để ở Chương 3.
% Văn phong: đoạn văn liên kết, câu chủ động có chủ ngữ thật ("đồ án", "Ethereum", "MetaMask"...; hạn chế xưng
% "em", "tác giả" chỉ dùng ở Kết chương), mục con dùng \subsection*,
% giữ thuật ngữ tiếng Anh (smart contract, token, gas, testnet...). Xem ghi chú đầu Chap1.tex.
% ============================================================

\chapter{CÔNG NGHỆ SỬ DỤNG}
\label{chap:congnghe}
\label{chap:cosolythuyet}
% ============================================================
% HÌNH VẼ CHƯƠNG 2 (TikZ) — gọi bằng macro trong Chap2.tex
%   \HinhChuoiKhoi  : cấu trúc liên kết các khối          (fig:chuoikhoi)
%   \HinhLuongDApp  : luồng tương tác của dApp             (fig:luongdapp)
% ============================================================

% ---------- Cấu trúc liên kết các khối ----------
\newcommand{\HinhChuoiKhoi}{%
\begin{figure}[htbp]
\centering
\begin{tikzpicture}[
    >=Stealth,
    truong/.style={draw, minimum width=3.3cm, minimum height=0.55cm,
                   font=\footnotesize, fill=white, inner sep=1pt},
    tieude/.style={truong, fill=gray!20, font=\footnotesize\bfseries}]
  \foreach \i/\x/\ten in {0/0/N-1, 1/4.6/N, 2/9.2/N+1} {
    \node[tieude] (h\i) at (\x,0) {Khối \ten};
    \node[truong, below=0pt of h\i] (p\i) {Mã băm khối trước};
    \node[truong, below=0pt of p\i] (t\i) {Dấu thời gian};
    \node[truong, below=0pt of t\i] (m\i) {Gốc Merkle};
    \node[truong, below=0pt of m\i, minimum height=1cm, align=center] (g\i)
        {Danh sách\\giao dịch};
    \node[draw, thick, rounded corners=2pt, inner sep=3pt, fit=(h\i)(g\i)] (b\i) {};
  }
  \draw[->, thick] (p1.west) -- (b0.east |- p1.west);
  \draw[->, thick] (p2.west) -- (b1.east |- p2.west);
  \node[font=\large] at (-2.35,-1.4) {$\cdots$};
  \node[font=\large] at (11.55,-1.4) {$\cdots$};
\end{tikzpicture}
\caption[Cấu trúc liên kết giữa các khối trong chuỗi khối]{Cấu trúc liên kết giữa các khối trong chuỗi khối \hinhtao[Claude]{20/10/2026}}
\label{fig:chuoikhoi}
\end{figure}}

% ---------- Luồng tương tác của ứng dụng phi tập trung ----------
\newcommand{\HinhLuongDApp}{%
\begin{figure}[htbp]
\centering
\begin{tikzpicture}[
    >=Stealth,
    hop/.style={draw, rounded corners=3pt, align=center, font=\footnotesize,
                text width=2.1cm, minimum height=1.1cm, fill=white},
    nhan/.style={font=\scriptsize, align=center}]
  \node[hop]                         (u)   at (0,0)     {Người dùng};
  \node[hop]                         (fe)  at (3.7,0)   {Giao diện web\\(Next.js,\\ethers.js)};
  \node[hop]                         (mm)  at (7.4,0)   {Ví\\MetaMask\\(giữ khóa riêng)};
  \node[hop, text width=2.6cm, minimum height=1.6cm, fill=gray!12]
                                     (net) at (11.3,0)  {Mạng Ethereum\\(Sepolia)\\nút RPC, EVM,\\hợp đồng thông minh};
  \node[hop]                         (be)  at (3.7,-3)  {Backend\\(Node.js, Express)};
  \node[hop]                         (db)  at (0,-3)    {CSDL off-chain\\(PostgreSQL)};

  \draw[->] (u) -- node[nhan, above] {thao tác} (fe);
  \draw[->] (fe) -- node[nhan, above] {yêu cầu ký\\(EIP-1193)} (mm);
  \draw[->] (mm) -- node[nhan, above] {giao dịch\\đã ký} (net);
  \draw[->, dashed] (fe.north) to[bend left=22]
      node[nhan, above] {đọc trạng thái on-chain (call)} (net.north);
  \draw[->] (net.south) |- node[nhan, pos=0.75, below] {sự kiện (Minted, Transfer, Retired\ldots)} (be.east);
  \draw[<->] (fe) -- node[nhan, right] {API REST} (be);
  \draw[<->] (be) -- node[nhan, above] {đọc/ghi} (db);
\end{tikzpicture}
\caption[Luồng tương tác giữa giao diện, ví MetaMask, mạng Ethereum và backend]{Luồng tương tác giữa giao diện, ví MetaMask, mạng Ethereum và backend \hinhtao[Claude]{20/10/2026}}
\label{fig:luongdapp}
\end{figure}}


% Tổng quan (Normal, không tiêu đề): chương trình bày các công nghệ mà đồ án dùng để hiện thực kiến trúc lai
% ở mục \ref{sec:dinhhuong}: blockchain và Ethereum, smart contract Solidity với OpenZeppelin và Hardhat, các chuẩn
% token ERC-721, ERC-20, cách ứng dụng web giao tiếp với blockchain qua MetaMask và ethers.js, và
% phần off-chain Node.js, PostgreSQL. Với mỗi công nghệ nêu yêu cầu nó phục vụ và lý do chọn.
% [VIẾT SAU] \canviet{Đoạn tổng quan chương}

% [VIẾT SAU] Toàn bộ các mục dưới đây chưa viết nội dung — bỏ dấu % để khôi phục.
% % ------------------------------------------------------------
% \section{Blockchain và nền tảng Ethereum}
% \label{sec:blockchain}
% \label{sec:ethereum}
%
% \subsection*{Blockchain}
% \label{subsec:cautrucblock}
% \label{subsec:dongthuan}
% % Ý cần có (ngắn gọn, chỉ phần đồ án dùng): khối, mã băm khối trước, sổ cái chỉ ghi thêm (ledger) ->
% % vì sao dữ liệu on-chain khó sửa, phục vụ yêu cầu toàn vẹn ở mục \ref{sec:phichucnang} và VĐ-2;
% % đồng thuận PoW/PoS, Ethereum chuyển sang PoS (The Merge 15/9/2022).
% % Trích: \cite{nakamoto2008bitcoin,zheng2017overview,ethereum2024pos}. Nhắc \autoref{fig:chuoikhoi} trước hình.
% % [VIẾT SAU] \canviet{Blockchain}
% %
% % \HinhChuoiKhoi
%
% \subsection*{Lựa chọn nền tảng blockchain}
% \label{subsec:phanloai}
% % Bảng so sánh Ethereum (public) với Hyperledger Fabric (permissioned, mà Vaccargiu dùng) và Chia
% % (CAD Trust dùng): mô hình truy cập, đồng thuận, ngôn ngữ smart contract, chuẩn token, ví, testnet, chi phí.
% % Lý do chọn Ethereum: phiếu giao nhiệm vụ; chuẩn ERC-721/ERC-20; MetaMask; testnet công khai để bên thứ ba
% % tự kiểm tra (VĐ-2, NC-07). Nhược điểm: phí gas, dữ liệu công khai -> xử lý ở yêu cầu quyền riêng tư.
% % [VIẾT SAU] \canviet{Lựa chọn nền tảng blockchain}
%
% \subsection*{Tài khoản, giao dịch và phí gas}
% \label{subsec:taikhoan}
% \label{subsec:evm}
% % Ý cần có: EOA và contract account; khóa riêng, địa chỉ ví; giao dịch ký số, nonce; EVM, gas, EIP-1559.
% % Đây là cơ sở cho phần đo gas ở Chương 3. Trích: \cite{buterin2014ethereum,wood2014yellow,eip1559}.
% % [VIẾT SAU] \canviet{Tài khoản, giao dịch và phí gas}
%
% \subsection*{Testnet Sepolia và Etherscan}
% \label{subsec:sepolia}
% % Ý cần có: mainnet và testnet; Sepolia, faucet; Etherscan để tra giao dịch, xác minh (verify) mã hợp đồng.
% % Lý do: phạm vi nguyên mẫu (mục \ref{sec:muctieu}), không tốn phí thật.
% % [VIẾT SAU] \canviet{Testnet Sepolia}
%
% % ------------------------------------------------------------
% \section{Smart contract và công cụ phát triển}
% \label{sec:smartcontract}
%
% \subsection*{Smart contract và ngôn ngữ Solidity}
% \label{subsec:khainiemsc}
% \label{subsec:solidity}
% % Ý cần có: định nghĩa (Szabo); tự thực thi, không sửa được mã sau triển khai; event và log (dùng cho
% % bộ đồng bộ sự kiện); cấu trúc contract (state variable, function, modifier, event), kế thừa.
% % Yêu cầu phục vụ: NC-03, NC-06 (ràng buộc sê-ri, không hủy hai lần thành ràng buộc kỹ thuật).
% % Trích: \cite{szabo1997formalizing,buterin2014ethereum,solidity2024docs}.
% % [VIẾT SAU] \canviet{Smart contract và Solidity}
%
% \subsection*{Thư viện OpenZeppelin và Hardhat}
% \label{subsec:openzeppelin}
% % OpenZeppelin 5: ERC721, ERC20, AccessControl (VERIFIER_ROLE...), Pausable, ReentrancyGuard -> NC-09 và
% % yêu cầu bảo mật. Hardhat: biên dịch, Hardhat Network, test Mocha/Chai, coverage, gas reporter.
% % Phương án thay thế: Foundry, Truffle (đã ngừng phát triển); lý do chọn Hardhat (JavaScript cùng ngôn ngữ backend).
% % Trích: \cite{openzeppelin2024contracts}.
% % [VIẾT SAU] \canviet{OpenZeppelin và Hardhat}
%
% \subsection*{Các lỗ hổng thường gặp của smart contract}
% \label{subsec:baomat}
% % Reentrancy, kiểm soát truy cập sai, tràn số (Solidity >= 0.8 tự kiểm tra), tx.origin. Mỗi lỗ hổng
% % 1–2 câu và cách phòng trong đồ án; công cụ quét Slither. Trích: \cite{atzei2017survey}.
% % [VIẾT SAU] \canviet{Lỗ hổng smart contract}
%
% % ------------------------------------------------------------
% \section{Các chuẩn token}
% \label{sec:token-standards}
%
% \subsection*{ERC-20 và ERC-721}
% \label{subsec:erc20}
% \label{subsec:erc721}
% % ERC-20: fungible token, totalSupply, balanceOf, transfer, approve, transferFrom, allowance \cite{eip20}.
% % ERC-721: NFT, ownerOf, safeTransferFrom, tokenURI, onERC721Received \cite{eip721}.
% % [VIẾT SAU] \canviet{ERC-20 và ERC-721}
%
% \subsection*{Lựa chọn chuẩn token cho tín chỉ các-bon}
% \label{subsec:chontoken}
% % Bảng so sánh ERC-20 / ERC-721 / ERC-1155 theo: giữ metadata riêng từng lô, chia nhỏ, thanh khoản,
% % hỗ trợ ví. Kết luận: mỗi lô là một NFT ERC-721 (giữ dự án, vintage, dãy sê-ri, NC-03), pool ERC-20
% % cho giao dịch từng phần (NC-05); kế thừa mô hình Batch NFT/TCO2/BCT của Toucan (mục \ref{sec:khaosatsobo}).
% % [VIẾT SAU] \canviet{Lựa chọn chuẩn token}
%
% % ------------------------------------------------------------
% \section{Ứng dụng web giao tiếp với blockchain}
% \label{sec:dapp}
%
% \subsection*{Kiến trúc dApp}
% \label{subsec:kientrucdapp}
% % Luồng một giao dịch từ giao diện tới smart contract (giải thích \autoref{fig:luongdapp}).
% % [VIẾT SAU] \canviet{Kiến trúc dApp}
% %
% % \HinhLuongDApp
%
% \subsection*{Ví MetaMask, ethers.js và Next.js}
% \label{subsec:metamask}
% \label{subsec:ethers}
% % MetaMask giữ khóa riêng, EIP-1193 (window.ethereum) \cite{eip1193}; ethers.js: Provider, Signer,
% % Contract, ABI, call và send \cite{ethers2024docs}; Next.js cho giao diện. Yêu cầu phục vụ: tính dễ dùng,
% % người dùng tự ký giao dịch (backend không giữ khóa). Phương án thay thế: web3.js, viem; React thuần.
% % [VIẾT SAU] \canviet{MetaMask, ethers.js và Next.js}
%
% % ------------------------------------------------------------
% \section{Công nghệ phía off-chain}
% \label{sec:offchain-congnghe}
% \label{sec:onoffchain-lythuyet}
% % Node.js + Express (API RESTful), PostgreSQL (dữ liệu theo CAD Trust V2.0, NC-10), JWT, xác thực bằng
% % chữ ký ví. Tiêu chí chọn dữ liệu on-chain/off-chain theo \cite{xu2017taxonomy}; neo mã băm keccak256 của hồ sơ;
% % IPFS \cite{benet2014ipfs} là hướng mở rộng (hiện lưu tệp cục bộ). Phương án thay thế: MongoDB; lý do chọn
% % CSDL quan hệ (mô hình CAD Trust có quan hệ rõ, giống UNDP dùng PostgreSQL).
% % [VIẾT SAU] \canviet{Công nghệ off-chain}

% ------------------------------------------------------------
% Kết chương (Normal, không tiêu đề): tóm tắt bộ công nghệ đã chọn (có thể kèm bảng tổng hợp
% Công nghệ | Phiên bản | Vai trò | Yêu cầu phục vụ) và dẫn sang Chương 3.
% [VIẾT SAU] \canviet{Kết chương 2}
